\documentclass[10pt,a4paper]{amsart}
\usepackage[margin=19mm]{geometry}
\usepackage{amsmath,amssymb,mathtools}
\usepackage[scr=rsfs,cal=euler]{mathalfa}
\usepackage{xfrac}
\usepackage[unicode,hidelinks,bookmarks=false]{hyperref}
\newcommand{\ie}{i.e.\ }
\newcommand{\cf}{cf.\ }
\newcommand{\resp}{respectively\ }
\newcommand{\Subsection}{\subsection}
\providecommand{\nobreakdash}{\nobreak\hbox{-}}
\setlength{\emergencystretch}{2em}
\multlinegap0pt
\usepackage{amssymb}
\usepackage{tikz}
\usepackage[shortlabels]{enumitem}
\newtheorem{proposition}{Proposition}
\newcommand\W{{\mathcal W}}
\newcommand\eps{{\varepsilon}}
\newcommand\supp{{\rm supp\,}}
\title{\bf Localisation of pseudospectra on discrete groups}
\author{Simon N. Chandler-Wilde, Marko Lindner, Christian Seifert}
\date{Source: arXiv:2607.29354v1 (CC BY 4.0)}
\begin{document}
\maketitle

\noindent\emph{Source and licence.} \bf Localisation of pseudospectra on discrete groups. Version arXiv:2607.29354v1 (CC BY 4.0), distributed by arXiv under the Creative Commons Attribution 4.0 International licence, http://creativecommons.org/licenses/by/4.0/, as recorded in the arXiv licence metadata for this exact version and shown on the abstract page https://arxiv.org/abs/2607.29354v1. Full licence notice: \texttt{licenses/arxiv-2607.29354-CC-BY-4.0.txt}.\par\smallskip

\noindent\textit{Editorial scope.} Counted unit: the article's proposition prop:tau1 (source0.tex lines 892--906). The article's complete original proof of it is delivered with it (source0.tex lines 907--976, one \texttt{proof} environment of 70 lines). No mathematics is written, repaired or re-derived: the statement and the proof are the article's own lines, in the article's own notation, and no retained line is substituted or re-flowed.  Retained uncounted as the source-local context the proof needs: lines 405--407; lines 851--853; lines 867--869; lines 873--875.  Nothing was omitted from any retained span, so the package carries no omission record.  No retained line contains a citation, so the wrapper carries no bibliography and no bibliography entry.  No retained line contains a figure environment, an image-inclusion command or drawing code, so no image is delivered and the package declares no asset dependency.  Editorial formatting only, in addition: an AMS \texttt{amsart} preamble that restates the article's own declarations and macros used by the retained lines, namely the environments \texttt{proposition} on separate counters, together with the standard packages those lines need.  The retained environments print in this document as Proposition 1 in order of appearance, whereas the article prints the same items with the numbering of its own full document.  The complete native source of the article as distributed in the arXiv e-print for this exact version is bundled unchanged as \texttt{sources/206476/source0.tex}.
\begin{equation} \label{eq:AW}
A\ =\ \sum_{j\in G}M_{b^{(j)}}V_j,
\end{equation}

\begin{equation} \label{eq:patchA_tau1}
\ A_k^{(\tau_1)}\ :=\ AP_k
\end{equation}

\begin{equation} \label{eq:patchx}
\tilde x_k\ :=\ M_{V_kw}x,\quad\text{so that}\quad (\tilde x_k)_i\ =\ w_{i-k}\cdot x_i,\quad\text{i.e.}\quad (\tilde x_k)_{k+m}\ =\ w_m\cdot x_{k+m}
\end{equation}

\begin{equation}\label{eq:lem-trick1}
\sum_{k\in G} \|\tilde x_k\|^2\, =\, \sum_{k\in G}\|M_{V_kw}x\|^2\, =\, \sum_{k\in G}\sum_{i\in G}\|w_{i-k}x_i\|^2\, =\, \sum_{i\in G}\|x_i\|^2\sum_{k\in G}|w_{i-k}|^2\, =\, \|x\|_2^2\,\|w\|_2^2,
\end{equation}

\begin{proposition} \label{prop:tau1}
Let $A\in \W(E)$, in which case $A$ has the representation \eqref{eq:AW},
for some sequence $(b^{(j)})_{j\in G} \subset \ell^\infty(G,L(Y))$ with $\|A\|_\W <\infty$, and 
suppose a finite set $W\subset G$ and  %a nonzero weight function
 $w\in \ell^2(G)$ with $w\neq 0$ and $\supp w\subset W$ 
 are given. Then, for every nonzero $x\in E$, there exists a $k\in G$ such that $\tilde x_k\ne 0$ and
\begin{equation} \label{eq:patchAx_tau1}
\frac{\|A_{k}\, \tilde x_{k}\|}{\| \tilde x_{k}\|}\ \le\ \frac{\|Ax\|}{\|x\|}\ +\ \eps(w,A),
\end{equation}
where
\begin{equation} \label{eq:pen_tau1}
\eps(w,A)\ :=\ \eps^{(\tau_1)}(w,A)\ :=\ \sum_{j\in G} \|b^{(j)}\|_\infty \frac{\|V_jw-w\|_2}{\|w\|_2}
\end{equation}
and where the local patches $A_k$ and $\tilde x_k$ are defined by \eqref{eq:patchA_tau1} and \eqref{eq:patchx} above.
\end{proposition}

\begin{proof}
Take a nonzero $x\in E$. Recalling the commutator notation $[A,B]:=AB-BA$ for bounded operators $A,B$, we have, for every $k\in G$,
\begin{eqnarray*}
A_k\tilde x_k &=& AP_kM_{V_kw}x \ =\ AM_{V_kw}x \ =\ M_{V_kw}Ax\ +\ [A,M_{V_kw}]x\\
&=& M_{V_kw}Ax\ +\ \left[\sum_{j\in G} M_{b^{(j)}}V_j\ ,\ M_{V_kw}\right]x
\ =\ M_{V_kw}Ax\ +\ \sum_{j\in G} [M_{b^{(j)}}V_j,M_{V_kw}]x,
\end{eqnarray*}
so that
\[
\|A_k\tilde x_k\| \ \le\ \|M_{V_kw}Ax\|\ +\ \sum_{j\in G} \|[M_{b^{(j)}}V_j,M_{V_kw}]x\|.
\]
%\|AM_{V_kw}x\|^2 \ \le\ \left(\|M_{V_kw}Ax\|\ +\ \sum_{j\in J} \|[M_{b^{(j)}}V_j,M_{V_kw}]x\|\right)^2\\
Applying squares, summation over all $k\in G$ and a square root, it follows that
\begin{eqnarray}
%\sum_{k\in G} \|A\tilde x_k\|^2 &=&
%\sum_{k\in G}\|AM_{V_kw}x\|^2 \ \le\ \sum_{k\in G}\left(\|M_{V_kw}Ax\|\ +\ \sum_{j\in J} \|[M_{b^{(j)}}V_j,M_{V_kw}]x\|\right)^2\\
\sqrt{\sum_{k\in G} \|A_k\tilde x_k\|^2} &\le&
%\sqrt{\sum_{k\in G}\|AM_{V_kw}x\|^2} \ \le\ 
\sqrt{\sum_{k\in G}\left(\underbrace{\|M_{V_kw}Ax\|}_{=:\ \xi_k}\ +\ \sum_{j\in G} \underbrace{\|[M_{b^{(j)}}V_j,M_{V_kw}]x\|}_{=:\ \eta_{j,k}}\right)^2}\nonumber\\ 
&=& \sqrt{\sum_{k\in G}(\xi_k+\sum_{j\in G} \eta_{j,k})^2}\ 
%=\ \|\xi+\sum_{j\in J}\eta_j\|
\ \stackrel{(M)}\le%\ \|\xi\|+\sum_{j\in J}\|\eta_j\|\ =
\ \sqrt{\sum_{k\in G} \xi_k^2}+\sum_{j\in J}\sqrt{\sum_{k\in G}\eta_{j,k}^2}\nonumber\\
&=& \sqrt{\sum_{k\in G}\|M_{V_kw}Ax\|^2}\ +\ \sum_{j\in G}\sqrt{\sum_{k\in G}\|[M_{b^{(j)}}V_j,M_{V_kw}]x\|^2} \label{eq:tag1}
\end{eqnarray}
with $(M)$ denoting an application of Minkowski's inequality.

Now we evaluate the right-hand side of \eqref{eq:tag1}. First, by \eqref{eq:lem-trick1},
\begin{equation} \label{eq:trick1}
\sqrt{\sum_{k\in G}\|M_{V_kw}Ax\|^2}=\sqrt{\|w\|_2^2\|Ax\|^2}=\|w\|_2\|Ax\|.%\le\|w\|(\nu(A)+\delta)$
\end{equation}
%
%Then we bound $\|[M_{b^{(j)}}V_j,M_{V_kw}]x\|^2$
Further, for $j,k\in G$,
\begin{eqnarray*}
 \|[M_{b^{(j)}}V_j,M_{V_kw}]x\|^2 &=& \sum_{i\in G}|(M_{b^{(j)}}V_jM_{V_kw}x)_i-(M_{V_kw}M_{b^{(j)}}V_jx)_i|^2\\
%&=& \sum_{i\in G}|b^{(j)}_i\cdot (M_{V_kw}x)_{i-j}-(V_kw)_i\cdot b^{(j)}_i\cdot (V_jx)_i|^2\\
&=& \sum_{i\in G}\|b^{(j)}_i \cdot w_{i-j-k}\cdot x_{i-j}-w_{i-k}\cdot b^{(j)}_i\cdot x_{i-j}\|^2\\
&\le& \|b^{(j)}\|_\infty^2\sum_{i\in G} |w_{i-j-k}-w_{i-k}|^2\|x_{i-j}\|^2.
%&\stackrel{m:=i-k}=& \|b^{(j)}\|_\infty^2\sum_{m\in G} |w_{m-j}-w_{m}|^2|x_{m+k-j}|^2,
\end{eqnarray*}
Substituting\footnote{Note that we use commutativity of the group operation when we turn $i-j-k$ into $i-k-j=m-j$.} $m:=i-k$, summing up over $k\in G$, and arguing as in \eqref{eq:lem-trick1},
\begin{eqnarray}
\sum_{k\in G} \|[M_{b^{(j)}}V_j,M_{V_kw}]x\|^2 &\le& \|b^{(j)}\|_\infty^2 \sum_{k\in G} \sum_{m\in G} |w_{m-j}-w_{m}|^2\|x_{m+k-j}\|^2 \nonumber\\
%&=& \|b^{(j)}\|_\infty^2 \sum_{m\in G} \sum_{k\in G} |w_{m-j}-w_{m}|^2|x_{m+k-j}|^2\nonumber\\
&=& \|b^{(j)}\|_\infty^2 \sum_{m\in G} |w_{m-j}-w_{m}|^2\sum_{k\in G} \|x_{m+k-j}\|^2\nonumber\\
&=& \|b^{(j)}\|_\infty^2 \|V_jw-w\|_2^2\|x\|_2^2.\label{eq:RHS2}
%\\
%&=& \|b^{(j)}\|_\infty^2 \|V_jw-w\|_2^2. 
\end{eqnarray}
  
So from \eqref{eq:tag1} and our bounds \eqref{eq:trick1} and \eqref{eq:RHS2} on its right-hand side we conclude
\begin{eqnarray}
\sqrt{\sum_{k\in G}\|A_k\tilde x_k\|^2} &\le& \|w\|_2\|Ax\| + \sum_{j\in J} \|b^{(j)}\|_\infty \|V_jw-w\|_2\|x\|\nonumber\\
&=& \underbrace{\left(\frac{\|Ax\|}{\|x\|} + \sum_{j\in J} \|b^{(j)}\|_\infty \frac{\|V_jw-w\|_2}{\|w\|_2}\right)}_{=:\,c}\|w\|_2\,\|x\|.\label{eq:tag2}
%\\
%\ =\ \left(\nu(A)+\delta + \sum_{j\in J} \|b^{(j)}\|_\infty \frac{\|V_jw-w\|_2}{\|w\|_2}\right) \|w\|_2.
\end{eqnarray}
By \eqref{eq:lem-trick1} again, 
\[
\sum_{k\in G}\|\tilde x_k||^2\ =\ \sum_{k\in G}\|M_{V_kw}x||^2\ =\ \|w\|_2^2\,\|x\|^2,
\]
so that \eqref{eq:tag2} turns into
\[
\sqrt{\sum_{k\in G}\|A_k\tilde x_k\|^2} \ \le\ c\sqrt{\sum_{k\in G}\|\tilde x_k||^2}
\qquad\text{and hence}\qquad
\sum_{k\in G}\|A_k\tilde x_k\|^2 \ \le\ c^2\sum_{k\in G}\|\tilde x_k||^2.
\]
We conclude that, with $c$ as defined in \eqref{eq:tag2}, either i) $\|A_k\tilde x_k\|^2\le c^2\|\tilde x_k\|^2$ for all $k\in G$, where $\tilde x_k\ne 0$ for at least one of them or ii) $\|A_k\tilde x_k\|^2> c^2\|\tilde x_k\|^2$ for some $k$ but $\|A_k\tilde x_k\|^2< c^2\|\tilde x_k\|^2$ for others. Also then $\tilde x_k\ne 0$ for the latter. It remains to take the square root and divide by $\|\tilde x_k\|$.
\end{proof}

\end{document}
